\section{Dual-Epoch Security}
\label{sec:dualepoch}

The attack of \cref{sec:transfer} exposes cohort H almost as much as cohort P. This section
explains why, states the security notion that separates the two cohorts, and gives the conditions
under which the separation can be made to hold. Three results delimit what any storage design can
achieve: a transfer lemma showing that a content-placed layout exposes the two cohorts equally, a
capacity result showing that bounded per-partition capacity leaks the earlier cohort's load, and a
separation showing that a self-writing agent leaks cohort H through its own retrievals against any
store whose operations depend on content.

\subsection{A Transfer Lemma}
\label{sec:transfer-lemma}

The empirical finding that cohort H is exposed nearly as much as cohort P is not an artifact of a
particular attack; it follows from the placement rule.

\begin{lemma}[Transfer]
\label{lem:transfer}
Let records be assigned to partitions by a function of content alone, let cohorts H and P be
drawn from the same distribution over content, and let the adversary observe the partition of
each stored record. For any partition-labeling $\ell$ that the adversary computes without
reference to cohort membership, the expected topic-recovery accuracy of $\ell$ on cohort H equals
its expected accuracy on cohort P.
\end{lemma}

\begin{proof}[Proof sketch]
Placement depends only on content, so a record's partition is independent of its cohort given its
content. The labeling $\ell$ maps partitions to topics without reading cohort membership, so the
predicted label of a record depends only on its content. The true topic is a function of content.
Hence the event that $\ell$ is correct on a record depends only on its content, whose distribution
is identical across cohorts; the expectations are therefore equal.
\end{proof}

The third hypothesis is easily overlooked. Applying a partition label to a cohort-H record
requires knowing which partition holds it: a growing layout discloses that in the snapshot, and
against a layout of fixed-size sealed regions the lemma says nothing at all. The write channel of
\cref{sec:eval-write} is therefore a statement about cohort P unless another channel supplies
cohort-H membership, and \cref{sec:eval-multisnapshot} exhibits one. The other hypotheses name the
ways the cohorts' exposure can differ: subject drift breaks identical distributions, which
\cref{sec:eval-drift} measures together with how much an adversary inducing with archival content
recovers. Dual-epoch security is thus valuable not because historical data is inherently safer,
but because a store can be built whose adversary-visible representation carries no partition
labels at all.

\subsection{The Security Game}

Dual-epoch security is defined with respect to a leakage function
$\mathcal{L} = (\mathcal{L}^{\mathrm{snap}}, \mathcal{L}^{\mathrm{op}})$, where
$\mathcal{L}^{\mathrm{snap}}$ describes the state exposed at $\tau$ and $\mathcal{L}^{\mathrm{op}}$
describes each post-$\tau$ operation. Making the leakage explicit lets the definition cover the
operations a real agent performs, searches and deletions included.

The leakage is stated over \emph{plaintext operations}, in the style usual for structured
encryption, rather than over the transcript the store happens to produce. Writing it the other way,
as "the physical regions the operation reads or writes", makes the definition circular: the
transcript is then permitted to reveal itself, and the theorem says little. Let $\pi$ be a secret
relabeling of partitions, sampled once and used consistently. Then
\begin{itemize}
\item $\mathsf{put}(x)$ reveals $\pi(\mathrm{part}(x))$;
\item $\mathsf{search}(q)$ reveals $\pi$ applied to the probe set of $q$, together with the fresh
random extras drawn on that call;
\item $\mathsf{delete}$ and $\mathsf{update}$ reveal nothing at the time of the call, and an
update additionally reveals the new record's relabeled partition at its next flush;
\item every operation reveals the time at which it occurred.
\end{itemize}
The simulator samples $\pi$ itself, which is what makes explicit that the keyed permutation buys
consistency of the relabeling and nothing more. \Cref{sec:eval-permutation} measures how much even
that is worth against an adversary that observes reads.

Let $\Pi$ be a dynamic encrypted memory. The game $\mathbf{G}^{\mathrm{H}}_{\Pi,\mathcal{A}}$ is
played between a challenger and an adversary $\mathcal{A}$.
\begin{enumerate}
\item $\mathcal{A}$ outputs two \emph{admissible} histories $\sigma_0, \sigma_1$ of equal length,
each a sequence of $\mathsf{put}$ operations. Admissibility is defined in \cref{sec:conditions}.
\item The challenger samples $b \leftarrow \{0,1\}$, initializes $\Pi$, executes $\sigma_b$, and
performs a flush, so that no record of $\sigma_b$ remains staged. This point is $\tau$. The
challenger gives $\mathcal{A}$ the snapshot leakage $\mathcal{L}^{\mathrm{snap}}$.
\item $\mathcal{A}$ then issues operations \emph{adaptively}: having seen the leakage of every
earlier operation, it chooses the next one, from $\mathsf{put}$, $\mathsf{delete}$,
$\mathsf{update}$, and $\mathsf{search}$, and may reference records written in $\sigma_b$. After
each operation the challenger returns $\mathcal{L}^{\mathrm{op}}$.
\item $\mathcal{A}$ outputs $b' \in \{0,1\}$.
\end{enumerate}
The advantage is
$\AdvH_{\Pi,\mathcal{A}}(\lambda) = \bigl|\Pr[b'=b] - \tfrac{1}{2}\bigr|$.

Adaptivity matters because the $\tau$-arrival adversary is inherently adaptive: it induces writes
in response to what it has already observed. The hybrid argument in the supplementary material is
unchanged by it, because each hybrid step replaces a primitive rather than fixing the operation
sequence.

\begin{definition}[$\mathcal{L}$-dual-epoch security]
\label{def:dualepoch}
$\Pi$ is $\mathcal{L}$-dual-epoch secure if for every probabilistic polynomial-time $\mathcal{A}$
there is a simulator $\mathcal{S}$ such that the view $\mathcal{A}$ receives in
$\mathbf{G}^{\mathrm{H}}_{\Pi,\mathcal{A}}$ is computationally indistinguishable from the output of
$\mathcal{S}$ given only $\mathcal{L}$ evaluated on the operations $\mathcal{A}$ issued, and
consequently $\AdvH_{\Pi,\mathcal{A}}(\lambda) \le \negl(\lambda)$.
\end{definition}

Because the simulator sees the post-$\tau$ leakage but nothing about $\sigma_b$, the definition
states that the entire view is independent of which history produced cohort H. The flush at $\tau$
matters: a record staged but unflushed at $\tau$ would be written in full view of the adversary.
Equal lengths matter too: the maintenance cursor's resting position counts pre-$\tau$ flushes, and
global occupancy bounds how many further records the store accepts, so unequal lengths would make
the definition unsatisfiable for reasons unrelated to placement (\cref{sec:mnemos-write} gives the
deployment measure for the cursor).

\Cref{def:dualepoch} concerns indistinguishability of histories, not topic recovery. The two are
connected by a standard reduction.

\begin{corollary}[Predicate hiding]
\label{cor:predicate}
If $\Pi$ is $\mathcal{L}$-dual-epoch secure, then for any predicate $\phi$ on cohort-H contents,
in particular the topic of any record, no adversary predicts $\phi$ from its view with advantage
over the best cohort-oblivious predictor exceeding $2\,\AdvH_{\Pi,\mathcal{A}}(\lambda)$, where the
advantage is averaged over the distribution from which the histories are drawn.
\end{corollary}

The argument is the usual equivalence between semantic security and indistinguishability, with
the averaging over histories that lets a $\phi$-predictor fix a distinguishing pair; the
supplementary material spells it out. The corollary must not be read backwards: a measured
advantage is a \emph{lower} bound on what one classifier extracts, so \cref{sec:eval-game} reports
the upper end of each interval as the ceiling for the attacker it ran, and claims nothing beyond
it.

\subsection{Sufficient Conditions}
\label{sec:conditions}

A history is \emph{admissible} if, under either $\sigma_0$ or $\sigma_1$, no partition receives
more than its capacity of live records at any point during the game. Tombstoned records that have
not yet been purged still occupy their slots, so admissibility counts them. Admissibility is a
restriction on the histories the adversary may submit, not on the scheme, and it is necessary: a
scheme cannot hide whether a record overflowed a partition if the adversary is free to choose
histories that differ precisely in whether an overflow occurs. \Cref{prop:capacity} and
\cref{sec:eval-refusal} address what happens when a deployment leaves the admissible regime, which
under realistic skew it does early.

\begin{theorem}
\label{thm:main}
Let $\Pi$ store records in $B$ partitions of fixed capacity $C$, in a single file of $B$ equal
regions, and let it satisfy the following conditions.
\begin{itemize}
\item \textbf{S1 (fixed size).} Every file the store maintains has a length determined by the public
parameters alone.
\item \textbf{S2 (sealed slots).} Every slot is encrypted with an IND\$-CPA authenticated cipher under a
fresh nonce; a free slot is uniformly random bytes; and any cleartext index within a slot is the
output of a pseudorandom function of a uniformly random identifier.
\item \textbf{S3 (data-independent updates).} Each flush rewrites exactly $w+1$ whole regions: one per
staged record, uniformly random regions for the remainder up to $w$, and one region chosen by a
counter that depends only on the flush index. A deletion writes nothing until the next flush. A
refused write performs the same rewrites an accepted one would. The manifest is rewritten every
flush regardless of content.
\item \textbf{S4 (independent placement).} The partition table is independent of cohort H, and every
submitted history is admissible.
\end{itemize}
Then $\Pi$ is $\mathcal{L}$-dual-epoch secure for the $\mathcal{L}$ above, with
$\mathcal{L}^{\mathrm{snap}}$ empty beyond the public parameters.
\end{theorem}

\begin{proof}[Proof sketch]
The proof is a sequence of hybrids over the game, given in full in the supplementary material.
$\mathbf{G}_0$ is the real game. $\mathbf{G}_1$ replaces every pseudorandom function with a random function; a
distinguisher breaks the PRF. $\mathbf{G}_2$ replaces every ciphertext, live or free, with a
uniformly random string of the same length; a distinguisher breaks IND\$-CPA, and S1 guarantees
the lengths match. In $\mathbf{G}_2$ the snapshot is a collection of fixed-length uniform strings
whose sizes are public parameters, so by S1 and S2 it is independent of $b$. Each post-$\tau$
operation writes a set of whole regions determined by S3, whose identities are the staged records'
relabeled partitions, uniform choices, and a content-independent counter, none of which depends on
which history produced the records already present. Under S4 the partition table is independent of
$b$, and admissibility ensures no write is refused under either history, so the refusal rule of S3
is never exercised and cannot distinguish. Hence the whole view is a function of $\mathcal{L}$ and
public randomness, which the simulator reproduces.
\end{proof}

Padding a flush with rewrites of uniformly random regions rather than with dummy records keeps the
write pattern data-independent without consuming capacity, so admissibility can be stated in terms
of real load alone.

Rewriting whole regions (S3) makes a write reveal which region changed and never which slot, and
it is natural to assume that a per-slot write would reveal how full the region was. It need not.
Whether the index of a written slot carries occupancy depends on how the slot was chosen, and a
uniformly random choice among the free slots makes whole-region rewrites unnecessary for placement.

\begin{lemma}[Slot-index independence]
\label{lem:slot}
Let each accepted record occupy a slot drawn uniformly from the free slots of its region, and let
nothing else in the store depend on slot labels within a region. For any admissible pair
$\sigma_0, \sigma_1$ and any region $r$, the sequence of slot indices at which post-$\tau$ records
are placed in $r$ has the same distribution under $b = 0$ and $b = 1$.
\end{lemma}

\begin{proof}
The placement rule is invariant under any permutation of the slot labels of $r$, and by hypothesis
nothing else distinguishes them, so the labels are exchangeable. Hence the set $O_b(r)$ of slots
cohort H occupies at $\tau$ is a uniformly random subset of size $L_b(r)$, whatever the history.
Post-$\tau$ records draw uniformly without replacement from its complement, so $k$ of them occupy a
uniformly random ordered $k$-subset of the complement of a uniformly random $L_b(r)$-subset, which is
a uniformly random ordered $k$-subset of all $C$ labels whenever $k \le C - L_b(r)$. Admissibility
guarantees that bound under both histories, so the distribution does not depend on $b$. By S2 the
snapshot does not reveal $O_b(r)$, so the adversary cannot condition on it.
\end{proof}

The lemma is exactly as strong as its hypothesis, which fails for the maintenance step: freeing
tombstoned cohort-H records slot by slot writes as many slots as were deleted there, which depends
on the history, so the purge must remain a whole-region rewrite. Record placement need not;
\cref{sec:eval-game} measures all four combinations of placement rule and write granularity, and
only the lowest free slot, written alone, distinguishes the histories.

A second boundary is set by what the game reveals. Every post-$\tau$ put is in
$\mathcal{L}^{\mathrm{op}}$, so the lemma says nothing about hiding \emph{whether} a write
happened, which Mnemos-M's clocked flush promises. There padding must look like placement, and slot
by slot it does not: a slot written twice with no deletion between was padding the first time, so
over enough flushes the adversary counts the puts. Per-slot placement is therefore available to the
partitioned design and not to the memory-resident one, where it would have saved writing $C$ slots
per record.

\subsection{Separations}

The conditions are not implied by the standard notions.

\begin{proposition}[Forward and backward privacy are insufficient]
\label{prop:fp}
There is a scheme that is forward and Type-II backward private but not dual-epoch secure.
\end{proposition}
A forward- and backward-private scheme may keep per-partition state, such as an update counter,
in its local snapshot. Counters sit behind pseudorandom handles, so the adversary sees a
\emph{multiset} of values rather than a vector indexed by partition; two histories with the same
per-partition counts in different partitions are therefore \emph{not} distinguished. Two histories
with different count \emph{profiles} are: writing $m$ records to one partition yields the multiset
$\{m\}$, while writing $m/2$ to each of two yields $\{m/2, m/2\}$, and no relabeling identifies
them. So $\AdvH$ is non-negligible while forward and backward privacy, properties of the operation
transcript, are unaffected. \Cref{sec:attack-counting} reports how much this state leaks on a
realistic corpus, where it is weak.

\begin{proposition}[Sealing and padding are each necessary]
\label{prop:seal}
Slot-level history independence with regions that grow with occupancy is not dual-epoch secure,
and padding regions to a fixed size without sealing free slots is not dual-epoch secure.
\end{proposition}
The first store leaks per-partition occupancy through region sizes even though within a region no
order is revealed; the second leaks occupancy through the boundary between written and unwritten
slots. \Cref{sec:eval-game} realizes both as ablations and measures the resulting advantage.

\subsection{Capacity and Confidentiality}
\label{sec:capacity-prop}

Admissibility is a real restriction, and a deployment cannot simply assume it.

\begin{proposition}[Capacity--confidentiality trade-off]
\label{prop:capacity}
Let $\Pi$ place records by content into partitions of capacity $C$, and let a write whose partition
is full be refused without performing the writes an accepted record would have performed. Then an
adversary that can induce writes on a topic and observe the operation leakage learns whether that
topic's live load has reached $C$. Since the adversary knows its own contribution to that load, it
learns the load contributed by cohort H. Dual-epoch security therefore requires either that $C$
exceed the worst-case per-topic load, or that a refusal be indistinguishable from an acceptance.
\end{proposition}

The proof is immediate: an accepted put reveals $\pi(\mathrm{part}(x))$ and a silently refused
one reveals nothing, and the distinguishing event is exactly "this topic's load is at capacity".
The second disjunct is the cheap one and is part of S3: the store owes the intended region its
rewrite whether or not the record was admitted. The first is expensive, since capacity must then be
provisioned for the most frequent topic; \cref{sec:eval-refusal} measures both. A design with only
a global capacity bound, such as \cref{sec:mnemos-m}, escapes the proposition, because global
occupancy is a function of the record count, which the game already equalizes.

\subsection{The Derived-Write Channel}
\label{sec:derived-write}

\Cref{def:dualepoch} treats the post-$\tau$ operations as chosen by the adversary. A self-writing
agent introduces a second source: it retrieves a cohort-H record and then, in the same turn, writes
a new record or issues a follow-up query about the same subject, so the operations performed after
$\tau$ depend on the contents of cohort H.

\begin{proposition}[Derived-write separation]
\label{prop:derived}
Let $\Pi$ be a store whose post-$\tau$ operation leakage $\mathcal{L}^{\mathrm{op}}$ reveals a
non-constant function of the content of the operations performed, or of whether an operation was
performed at all. Then there is a feedback agent, whose post-$\tau$ operations depend on retrieved
cohort-H records, against which $\Pi$ is not dual-epoch secure.
\end{proposition}

\begin{proof}
Consider an agent that, after $\tau$, searches for a fixed query $q$ and writes a record whose text
is its top result. Choose $\sigma_0$ and $\sigma_1$ so that the top result for $q$ lies in
different partitions under the two histories, which is possible because the histories are the
adversary's to choose. The post-$\tau$ write then touches a region that differs between $b=0$ and
$b=1$, and by hypothesis $\mathcal{L}^{\mathrm{op}}$ reveals a non-constant function of it. The
adversary outputs the corresponding bit and wins with non-negligible advantage. The same argument
runs on the occurrence of a write rather than its content: if the agent deduplicates, so that it
writes only when the retrieved fact is absent, then whether any write occurred is itself the
distinguishing bit.
\end{proof}

The proposition does not say that no storage design closes this channel; its hypothesis excludes
stores whose operation leakage is content-independent. Oblivious RAM is one. So is a store with
content-oblivious placement, in-memory search, and a clocked flush: the regions written are uniform
whatever the agent retrieved, and a flush occurs every $\Delta$ whether or not the agent wrote.
\Cref{sec:mnemos-m} is such a store. This is what separates self-writing memory from classical
searchable encryption, whose queries are external to the stored data: closing the channel needs
content-oblivious access \emph{and} a content-independent operation schedule, and
\cref{sec:eval-dedup} measures the resulting attack, deduplication membership inference, on both
sides of that line.
